\label{sec:human_eval}


\begin{table}
\centering
\resizebox{\columnwidth}{!}{%
\begin{tabular}{lccccc} 
\toprule
                & \multicolumn{2}{c}{oRAG}    & \multicolumn{2}{c}{STORM}       & \multirow{2}{*}{$p$-value}  \\
                & Avg.          & $\geq 4$ Rates & Av.g.         & $\geq 4$ Rates      &                           \\ 
\midrule
Interest Level  & 3.63          & 57.5\%     & \textbf{4.03} & \textbf{70.0\%}          & 0.077                     \\
Organization    & 3.25          & 45.0\%     & \textbf{4.00} & \textbf{70.0\%}          & 0.005                     \\
Relevance       & 3.93          & 62.5\%     & \textbf{4.15} & \textbf{65.0\%}          & 0.347                     \\
Coverage        & 3.58          & 57.5\%     & \textbf{4.00} & \textbf{67.5\%}          & 0.084                     \\
Verifiability   & \textbf{3.85} & 67.5\%     & 3.80          & 67.5\%          & 0.843                     \\ 
\midrule
\#Preferred & \multicolumn{2}{c}{14}     & \multicolumn{2}{c}{\textbf{26}} &                           \\
\bottomrule
\end{tabular}}
\caption{Human evaluation results on 20 pairs of articles generated by \system and \textit{oRAG}. Each pair of articles is evaluated by two Wikipedia editors.
The ratings are given on a scale between 1 and 7, with values $\geq4$ indicating good quality (see~\reftab{table:human_eval_rubric}). We conduct paired $t$-test and report the $p$-value.}
\label{table:human_eval}
\vspace{-0.7em}
\end{table}


To better understand the strengths and weaknesses of \system, we conduct human evaluation by collaborating with 10 experienced Wikipedia editors who have made at least 500 edits on Wikipedia and have more than 1 year of experience. We randomly sample 20 topics from our dataset and evaluate the articles generated by our method and \textit{oRAG}, the best baseline according to the automatic evaluation. Each pair of articles is assigned to 2 editors.

We request editors to judge each article from the same five aspects defined in~\refsec{sec:automatic_metrics}, but using a 1 to 7 scale for more fine-grained evaluation. While our automatic evaluation uses citation quality as a proxy to evaluate \textit{Verifiability}, we stick to the Wikipedia standard of ``verifiable with no original research'' in human evaluation. 
Besides rating the articles, editors are asked to provide open-ended feedback and pairwise preference. After the evaluation finishes, they are further requested to compare an article produced by our method, which they have just reviewed, with its human-written counterpart, and report their perceived usefulness of \system using a 1-5 Likert scale. More human evaluation details are included in Appendix~\ref{appendix:human_eval}. \reftab{table:human_eval} presents the rating and pairwise comparison results.\footnote{For the 1-7 scale rating results on each criterion, we calculate the Krippendorff's Alpha to measure the inter annotator agreement (IAA), and the results are as follows: \textit{Interest Level} (0.349), \textit{Organization} (0.221), \textit{Relevance} (0.256), \textit{Coverage} (0.346), \textit{Verifiability} (0.388).}

\noindent
\textbf{Articles produced by \system exhibit greater breadth and depth than oRAG outputs.} In accord with the finding in %
\refsec{sec:main_result}, editors judge articles produced by \system as more interesting, organized, and having broader coverage compared to %
\textit{oRAG} outputs. Specifically, 25\% more articles produced by \system are considered organized (\textit{Organization} rating $\geq 4$), and 10\% more are deemed to have good coverage (\textit{Coverage} rating $\geq 4$). Even in comparison with human-written articles, one editor praises our result as providing ``a bit more background information'' and another notes that ``I found that the AI articles had more depth compared to the Wikipedia articles''. \system also outperforms the best baseline in pairwise comparison.

\noindent
\textbf{More information in $|\mathcal{R}|$ poses challenges beyond factual hallucination.} We examine 14 pairwise comparison responses where editors prefer oRAG outputs over \system. Excluding 3 cases where pairwise preferences do not align with their ratings, editors assign lower \textit{Verifiability} scores to articles from our approach in over 50\% of the cases. Through analyzing the articles and editors' free-form feedback, 
we discover that \textit{low Verifiability scores stem from red herring fallacy or overspeculation issues}. These arise when the generated articles introduce unverifiable connections between different pieces of information in $|\mathcal{R}|$ or between the information and the topic (examples included in~\reftab{table:comment_summary}). Compared to the widely discussed factual hallucination~\citep{shuster-etal-2021-retrieval-augmentation,huang2023survey}, addressing such verifiability issues is more nuanced, surpassing basic fact-checking~\citep{min-etal-2023-factscore}. %


\begin{figure}[t]
    \centering 
    \resizebox{\columnwidth}{!}{%
    \includegraphics{img/survey.pdf}
    }
    \caption{Survey results of the perceived usefulness of \system ($n=10$).
    }
    \label{Fig.survey}
\vspace{-0.5em}
\end{figure}

\noindent
\textbf{Generated articles trail behind well-revised human works.}
While \system outperforms the oRAG baseline, editors comment that the generated articles are \textit{less informative than actual Wikipedia pages}. Another major issue identified is \textit{the transfer of bias and tone from Internet sources to the generated article}, with 7 out of 10 editors mentioning that the STORM-generated articles sound ``emotional'' or ``unneutral''. More analysis is discussed in Appendix~\ref{appendix:error_analysis}. This feedback suggests that reducing the retrieval bias in the pre-writing stage is a worthwhile direction for future work.

\noindent
\textbf{Generated articles %
are a good starting point.} As shown in~\reffig{Fig.survey}, editors are unanimous in agreeing that \system can aid them in their pre-writing stage.
It is gratifying to know that the tool is helpful to experienced editors.
80\% of the editors think that \system can help them edit a Wikipedia article for a new topic. More reservation is expressed to the usefulness of \system for the Wikipedia community at large; nonetheless, 70\% of the editors think it is useful, with only 10\% disagreeing. 




